% !TEX root = ../main.tex
\chapter{Differential structures: the pivotal concept of tangent vector spaces}\label{ch:l09}
\begin{paracol}{2}

\begin{sources}
\begin{tabularx}{\linewidth}{@{}l l L@{}}
\toprule
\textbf{Video} & \textbf{Length} & \textbf{Content}\\
\midrule
\seg{9}{1}{Lecture 9: Differential structures: the pivotal concept of tangent vector spaces} & 1:44:15 & The vector space $C^\infty(M)$; tangent vectors as directional derivative operators along curves;
the tangent space $T_pM$ and the proof that addition and scalar multiplication close; an aside on algebras, derivations, Lie algebras and the Poisson bracket; the theorem $\dim T_pM=\dim M$ with the chart-induced basis
$\partial/\partial x^a$; the difference between $\partial_a$ and $\partial/\partial x^a$; legality of acting on coordinate functions; components, change of chart and why there are no position vectors.\\
\bottomrule
\end{tabularx}

\smallskip
\textbf{Transcript} \file{materials/transcripts/C09-V01.txt} (YouTube automatic captions).
\end{sources}

\begin{reading}[Where this lecture sits in the other notes]
\emph{Gravity and Light}, Lecture~5 (Tangent spaces); \emph{Theoretische Mechanik}, Lecture~4. Both construct the tangent space of a smooth manifold; the aside on algebras and derivations, the careful distinction between the partial
derivative $\partial_a$ on $\R^D$ and the tangent vector $\partial/\partial x^a$, and the remarks on position vectors are, as far as the titles say, emphasised in this course.
(From the lecture titles and the plan of the course; not checked in detail.)
\end{reading}

% ==================================================================
\section{The vector space of smooth functions}\label{sec:tv-cinf}

Today the tangent spaces to a manifold: the definitions made here are ``probably the single most important definitions in differential geometry'', that of a tangent vector to a manifold and everything built from it \ts{9}{1}{0:09}. Having reviewed vector spaces last time
(\cref{ch:l08}), including infinite dimensional ones, one makes use of that now: the tangent space will be a finite dimensional vector space if the manifold is finite dimensional, with the same dimension as the manifold has as a topological manifold. But first an \emph{infinite
dimensional} vector space is constructed from a smooth manifold, an object that will return again and again \ts{9}{1}{0:30} \ts{9}{1}{0:51}.

Let $M$ be a smooth manifold. Throughout this chapter ``smooth manifold'' comes with a topology and a smooth atlas (\cref{def:df-manifold}); this is suppressed in the notation, and made explicit again only when several manifolds with several topologies and atlases are compared. ``From now on we are grown up'' \ts{9}{1}{1:11} \ts{9}{1}{1:32}.

\begin{definition}[New concepts: the vector space $C^\infty(M)$]\label{def:tv-cinf}
Let $M$ be a smooth manifold. The set
\[
  C^\infty(M):=\{f\colon M\to\R\mid f\text{ smooth}\}
\]
is a vector space over $\R$ (the manifolds are real) with the \textbf{pointwise} operations: for $f,g\in C^\infty(M)$ and $\lambda\in\R$
\[
  (f+g)(p):=f(p)+g(p),\qquad (\lambda\cdot f)(p):=\lambda\,f(p)\qquad\text{for all }p\in M,
\]
where on the right-hand sides are the addition and the multiplication of the real numbers, since $f$ and $g$ send points to real numbers \ts{9}{1}{1:52} \ts{9}{1}{2:13} \ts{9}{1}{2:54} \ts{9}{1}{3:36}.
\end{definition}
That this is a vector space is quickly checked (``pretty obvious''), and it is an infinite dimensional vector space which will play an important role \ts{9}{1}{3:56} \ts{9}{1}{4:17}.
(Smoothness of $f$ is meant in the sense of \cref{def:df-map}, with $\R$ carrying its single standard chart.)

% ==================================================================
\section{Tangent vectors as directional derivative operators}\label{sec:tv-vector}

The central concept follows \ts{9}{1}{4:17}. A curve is a map from the real line; to avoid worrying about intervals (``one could be very picky and say it is an interval from $a$ to $b$ and then one must worry about the domains''; one can always reparametrise) one is a little cavalier:
for the purposes here a curve has the full real line as parameter range; more precision will be added when it matters \ts{9}{1}{4:39} \ts{9}{1}{4:59} \ts{9}{1}{5:20}.

\begin{definition}[New concepts: tangent vector, directional derivative operator]\label{def:tv-vector}
Let $\gamma\colon\R\to M$ be a smooth curve through a point $p\in M$. Without loss of generality (shift the parameter) $\gamma(0)=p$; this could be dropped but makes the formulae simpler \ts{9}{1}{5:40} \ts{9}{1}{6:02} \ts{9}{1}{6:23}.
The \textbf{directional derivative operator at $p$ along $\gamma$} is the linear map
\[
  X_{\gamma,p}\colon C^\infty(M)\longrightarrow\R,\qquad f\longmapsto X_{\gamma,p}(f):=(f\circ\gamma)'(0).
\]
Here $f\circ\gamma\colon\R\to\R$ is a map from the real line to itself, for which the derivative is known (the prime), evaluated at the parameter value $0$ at which $\gamma$ passes through $p$: the result is a real number. The target $\R$ is regarded as a real vector space
over itself, $\R^1$, so that linearity makes sense on both sides. It is linear essentially because it is built with a derivative \ts{9}{1}{6:43} \ts{9}{1}{7:26} \ts{9}{1}{8:07} \ts{9}{1}{8:28} \ts{9}{1}{9:09} \ts{9}{1}{9:50}.
In differential geometry $X_{\gamma,p}$ is called the \textbf{tangent vector to the curve $\gamma$ at the point $p\in M$} \ts{9}{1}{10:11} \ts{9}{1}{10:32}.
\end{definition}
(The name \emph{tangent}: from Latin \emph{tangere}, to touch; not said in class.)
The terminology should not mislead one into thinking that this is anything like a collection of numbers; ``finally it will be'', but at this point the tangent vector, formally and abstractly, \emph{is} a directional derivative operator: tangent vectors are operators, derivative operators
acting on smooth functions. That is how the concept of a vector survives in differential geometry \ts{9}{1}{10:53} \ts{9}{1}{11:14}.

\begin{intuition}[Direction and derivative are married]\label{int:tv-married}
From undergraduate physics one has a direction $\vec X$ and the $\nabla$ operator, and the dot product $\vec X\cdot\nabla$ gives a directional derivative: in vector calculus one has an independent notion of direction and an independent notion of derivative and combines them. In
differential geometry the two concepts are married and \emph{cannot be separated}: there is only the combination. This is the key point in the technical, hence the practical, understanding of tangent vectors, and the consequences will be worked out today \ts{9}{1}{11:14} \ts{9}{1}{11:35} \ts{9}{1}{11:55} \ts{9}{1}{12:15}.
\end{intuition}

\begin{intuition}[$X_{\gamma,p}$ is the velocity of $\gamma$ at $p$]\label{int:tv-velocity}
A precise intuition: $X_{\gamma,p}$ is the \emph{velocity} of the curve $\gamma$ at $p$. Recall that tangent spaces are to be constructed not by referring to an outside space in which the manifold is embedded (\cref{int:df-embedding}), but intrinsically: given a smooth
manifold and a curve through a point, with no reference to the bulk around it, one of the tangent vectors at the point is the velocity of that curve, which is a vector, by definition, namely $X_{\gamma,p}$ \ts{9}{1}{12:38} \ts{9}{1}{12:58} \ts{9}{1}{13:20} \ts{9}{1}{13:40}.
\end{intuition}
\emph{Consistency check.} Let $\delta(\lambda):=\gamma(2\lambda)$, the curve that runs through the same points as $\gamma$ at double the speed (also $\delta(0)=p$). Acting on an arbitrary smooth $f$, the chain rule gives an inner derivative:
\[
  X_{\delta,p}(f)=(f\circ\delta)'(0)=2\,(f\circ\gamma)'(0)=2\,X_{\gamma,p}(f),
\]
so $X_{\delta,p}=2X_{\gamma,p}$, and any other factor (three times, twenty times $\pi$) pulls out in the same way. Running through the curve at double speed gives a vector twice as long: this is how a velocity must scale. This is only an intuition; the definition is the technical one above
\ts{9}{1}{14:01} \ts{9}{1}{14:22} \ts{9}{1}{14:43} \ts{9}{1}{15:03} \ts{9}{1}{15:23} \ts{9}{1}{15:44}.

% ==================================================================
\section{The tangent space $T_pM$ and why it is closed}\label{sec:tv-space}

\begin{definition}[New concepts: tangent space]\label{def:tv-space}
The \textbf{tangent vector space to $M$ at $p\in M$}, $T_pM$, has as underlying set
\[
  T_pM:=\{X_{\gamma,p}\mid\gamma\text{ a smooth curve through }p\},
\]
the set of \emph{all} tangent vectors at $p$ to all smooth curves through $p$. It is equipped with two operations $\oplus$ and $\odot$ (called so for the moment, since it is not yet clear that they close on the set), defined \textbf{pointwise} as maps into $\R$: for
$X_{\gamma,p},X_{\delta,p}\in T_pM$, $\lambda\in\R$ and $f\in C^\infty(M)$
\[
  (X_{\gamma,p}\oplus X_{\delta,p})(f):=X_{\gamma,p}(f)+X_{\delta,p}(f),\qquad(\lambda\odot X_{\gamma,p})(f):=\lambda\,X_{\gamma,p}(f),
\]
using on the right the addition and multiplication in $\R$ \ts{9}{1}{16:08} \ts{9}{1}{16:31} \ts{9}{1}{18:37} \ts{9}{1}{19:41} \ts{9}{1}{20:24} \ts{9}{1}{21:05}.
\end{definition}
The intuition: pick $p$ and a curve; look also at the same curve run through at double speed, at $\mu$ times the speed, backwards (a tangent vector in the opposite direction); then take any other smooth curve through $p$ at any speed (technically a different map
$\R\to M$): the collection of all the corresponding directional derivative operators is the set underlying the tangent space. ``It's very simple'' \ts{9}{1}{16:51} \ts{9}{1}{17:12} \ts{9}{1}{17:33} \ts{9}{1}{17:54}.

\paragraph{The question of closure.} The sum of two linear maps $C^\infty(M)\to\R$ is again a linear map $C^\infty(M)\to\R$; that is not in doubt. But an element of $T_pM$ is not just any linear map: it is one constructed in the specific way $f\mapsto(f\circ\sigma)'(0)$ from a curve.
So one must find, for the sum, a curve $\sigma$ with
\[
  X_{\gamma,p}(f)+X_{\delta,p}(f)=X_{\sigma,p}(f)\qquad\text{for all }f\in C^\infty(M);
\]
only then can one be confident that the addition closes on $T_pM$. The same for the multiplication by scalars \ts{9}{1}{21:46} \ts{9}{1}{22:06} \ts{9}{1}{22:27} \ts{9}{1}{22:48} \ts{9}{1}{23:08}.

\begin{proposition}[$\oplus$ and $\odot$ close on $T_pM$]\label{prop:tv-closure}
For $X_{\gamma,p},X_{\delta,p}\in T_pM$ and $\lambda\in\R$ the maps $X_{\gamma,p}\oplus X_{\delta,p}$ and $\lambda\odot X_{\gamma,p}$ are again elements of $T_pM$; hence $(T_pM,\oplus,\odot)$ is a real vector space \ts{9}{1}{23:30} \ts{9}{1}{23:52}.
\end{proposition}
\begin{proof}
\emph{The sum.} Employ a chart $(U,x)$ with $p\in U$ (the curves may be restricted to $\gamma^{-1}(U)$; ``let us not be so picky'') and define
\[
  \sigma(\lambda):=x^{-1}\Bigl(\bigl(x\circ\gamma\bigr)(\lambda)+\bigl(x\circ\delta\bigr)(\lambda)-x(p)\Bigr).
\]
This looks funny; first see whether it makes sense. The maps $x\circ\gamma$ and $x\circ\delta$ are the chart representations of the two curves, maps from (a part of) $\R$ to $\R^D$, $D=\dim M$; the third term is the constant map $\lambda\mapsto x(p)$. As maps into $\R^D$ they can be added
(addition in $\R^D$), and the term $-x(p)$ ensures that at $\lambda=0$ the resulting curve passes through $x(p)$: indeed
$\sigma(0)=x^{-1}(x(p)+x(p)-x(p))=p$ \ts{9}{1}{24:13} \ts{9}{1}{24:59} \ts{9}{1}{25:21} \ts{9}{1}{26:02} \ts{9}{1}{27:03} \ts{9}{1}{29:07}.
The picture: one wants a curve in the manifold whose tangent vector is the sum, but curves on the manifold cannot be added, since the manifold has not enough structure; so one goes to the chart, where the chart images of the two curves through $x(p)$ can be added, and lifts the result back up with $x^{-1}$
\ts{9}{1}{27:44} \ts{9}{1}{28:05} \ts{9}{1}{28:25} \ts{9}{1}{28:46}.

This seems to be ill-defined: the same curves could be represented in very different charts, and the sum curve in one chart looks very different from the sum in another. If one defined the sum of the curves as such it would be wildly chart dependent; the curve $\sigma$ certainly is. What saves the day is
that one is not interested in the whole sum curve but only in its derivative at $0$, and the derivative takes care of all that: the sum of the tangent vectors is well defined \ts{9}{1}{29:28} \ts{9}{1}{29:48} \ts{9}{1}{30:09} \ts{9}{1}{30:30}.

Write $\hat f:=f\circ x^{-1}\colon\R^D\to\R$ for the chart representative of $f$. Then
\[\begin{aligned}
X_{\sigma,p}(f)&=(f\circ\sigma)'(0)\\
&=\bigl(\hat f\circ(x\circ\gamma+x\circ\delta-x(p))\bigr)'(0)\\
&=\partial_a\hat f(x(p))\,\bigl[(x^a\circ\gamma)'(0)\\
&\qquad+(x^a\circ\delta)'(0)\bigr]\\
&=(\hat f\circ x\circ\gamma)'(0)+(\hat f\circ x\circ\delta)'(0)\\
&=X_{\gamma,p}(f)+X_{\delta,p}(f).
\end{aligned}\]
The third line is the multi-dimensional chain rule ($\hat f$ is the chart representative of $f$), applied to $x\circ\gamma+x\circ\delta-x(p)\colon\R\to\R^D$; the derivative of the constant is zero. The point at which $\partial_a(f\circ x^{-1})$ is evaluated is $x(\sigma(0))=x(p)$. The fifth line
uses the chain rule backwards, together with $x^{-1}\circ x=\mathrm{id}$ and the associativity of the composition. So $X_{\gamma,p}\oplus X_{\delta,p}=X_{\sigma,p}\in T_pM$: ``we pull this stunt of going to a chart and constructing a sum of curves, something terribly ill-defined; what is not ill-defined is its first derivative,
independent of the chart because the chart map cancels out again. That's the trick'' \ts{9}{1}{30:51} \ts{9}{1}{31:53} \ts{9}{1}{33:18} \ts{9}{1}{34:20} \ts{9}{1}{36:06} \ts{9}{1}{37:07}.

\emph{The scalar multiplication} is much simpler: one does the same with $\lambda$ times $x\circ\gamma$ and ``does not need all the rest'' \ts{9}{1}{37:07} \ts{9}{1}{37:28}.
\end{proof}
\begin{coursenote}[Details left out in class]
(i) The curve $\sigma$ is defined only for those $\lambda$ for which $x\circ\gamma(\lambda)+x\circ\delta(\lambda)-x(p)$ lies in $x(U)$; this is a neighbourhood of $0$, which is all that the derivative at $0$ needs (and one may reparametrise to the full line). (ii) For the scalar multiplication the cleanest curve is the
reparametrisation $\tau(\lambda):=\gamma(\mu\lambda)$ with $\mu$ the scalar: $X_{\tau,p}(f)=\mu\,(f\circ\gamma)'(0)$, exactly as in the doubling computation above (\cref{int:tv-velocity}); the lecture's version is $x^{-1}\bigl(x(p)+\mu\,(x\circ\gamma-x(p))\bigr)$,
which gives the same vector. (iii) The zero vector is the tangent vector of the constant curve. Not said in class.
\end{coursenote}
Looking back at the definition of $T_pM$ one sees that it is a very simple one, once the tangent vectors are accepted to be these linear maps: addition and scalar multiplication are those of maps, and one has shown that they close. Now that one knows, one never needs to worry again:
for any sum or scalar multiple there is an underlying curve. Note also that there are many curves producing the same tangent vector, for instance different curves through $p$ with the same speed give rise, although they are different curves, to the same tangent vector; it suffices to know that some curve exists
\ts{9}{1}{37:49} \ts{9}{1}{38:10} \ts{9}{1}{45:53}.

% ==================================================================
\section{Aside: algebras and derivations}\label{sec:tv-derivation}

Before looking at the properties of the tangent vectors, a little aside on \emph{algebras and derivations} (a technical term; not the same as derivatives) \ts{9}{1}{38:10} \ts{9}{1}{38:32}.

\begin{definition}[New concepts: algebra]\label{def:tv-algebra}
A vector space $(V,+,\cdot)$ over a field $K$ equipped with a \textbf{product}, a bilinear map
\[
  \bullet\colon V\times V\to V,
\]
is called an \textbf{algebra}; so it is a quadruple, where besides the scalar multiplication with respect to the underlying field one has a product between the elements of the vector space itself \ts{9}{1}{38:53} \ts{9}{1}{39:14} \ts{9}{1}{39:36} \ts{9}{1}{39:57}.
\end{definition}
The product can take various forms (``we will meet several in the future''); in particular, the \emph{Lie product} will be an antisymmetric one; no symmetry is assumed, only bilinearity. An example from high school: the cross product on $\R^3$ makes $\R^3$ an algebra (``you can't do this on $\R^4$'')
\ts{9}{1}{40:18} \ts{9}{1}{40:39}.
\begin{coursenote}[On the cross product in $\R^4$]
The lecture says that one cannot do this on $\R^4$. A bilinear antisymmetric product certainly exists on every vector space; what exists only in dimensions $3$ and $7$ is a cross product with the properties of the one on $\R^3$ (orthogonal to both factors, norm $|a||b|\sin\theta$). I cannot confirm that
this is what is meant; [?]. Not said in class.
\end{coursenote}
\begin{example}[$C^\infty(M)$ is an algebra]\label{ex:tv-cinf-alg}
On $C^\infty(M)$ define $\bullet\colon C^\infty(M)\times C^\infty(M)\to C^\infty(M)$ pointwise by $(f\bullet g)(p):=f(p)\,g(p)$, the multiplication of the real numbers. This is the multiplication of two functions; whereas the scalar multiplication is only the multiplication with a fixed constant: ``it's structurally totally
different, very important to realise this''. So $C^\infty(M)$ is not only a vector space but, with this product, an $\R$ \emph{algebra}. It seems formal overkill, since one always knew how to multiply two functions, but the point is to realise the structural level of the whole thing
\ts{9}{1}{40:39} \ts{9}{1}{41:00} \ts{9}{1}{41:21} \ts{9}{1}{41:41} \ts{9}{1}{42:02}.
\end{example}

\begin{definition}[New concepts: derivation]\label{def:tv-derivation}
Let $(A,+,\cdot,\bullet_A)$ and $(B,+,\cdot,\bullet_B)$ be algebras. A \textbf{derivation} is a linear map $D\colon A\to B$ which in addition satisfies the \textbf{Leibniz rule}
\[
  D(f\bullet_A g)=D(f)\bullet_B g+f\bullet_B D(g)\qquad\text{for all }f,g\in A.
\]
Linearity refers to the vector space underlying the algebras; the Leibniz rule is the compatibility with the products \ts{9}{1}{42:22} \ts{9}{1}{42:46} \ts{9}{1}{43:07} \ts{9}{1}{43:49} \ts{9}{1}{44:09} \ts{9}{1}{46:57}.
\end{definition}
(In class the definition was first written for a single algebra and then generalised, after some hesitation, to a map between two algebras.) The Leibniz rule terribly resembles the product rule of differentiation, and that is no accident: ordinary differentiation is a derivation, linear (the sum rule) and satisfying the Leibniz rule (the
product rule); the abstract version works in far more situations than the one used here \ts{9}{1}{44:30} \ts{9}{1}{44:51} \ts{9}{1}{45:11}.

\begin{example}[Tangent vectors are derivations]\label{ex:tv-der-tangent}
Take $A=C^\infty(M)$ with the pointwise product and $B=\R$ as an algebra over itself (all products are the multiplication in $\R$). Every $X\in T_pM$ is a derivation $C^\infty(M)\to\R$:
\[
  X(f\bullet g)=X(f)\,g(p)+f(p)\,X(g).
\]
One does not need to say to which curve and at which point $X$ corresponds: being in $T_pM$ it comes from some curve through $p$. This is quickly checked and goes to the problem sheet (\cref{ex:tv-leibniz}). So a tangent vector at $p$ is a particular example of a derivation
\ts{9}{1}{45:11} \ts{9}{1}{45:32} \ts{9}{1}{46:57} \ts{9}{1}{47:18} \ts{9}{1}{47:38} \ts{9}{1}{47:58}.
\end{example}
\begin{exercise}[Tangent vectors satisfy the Leibniz rule]\label{ex:tv-leibniz}
Show that $X_{\gamma,p}(f\bullet g)=X_{\gamma,p}(f)\,g(p)+f(p)\,X_{\gamma,p}(g)$ using the product rule for $(f\circ\gamma)\cdot(g\circ\gamma)\colon\R\to\R$. (Problem sheet.)
\end{exercise}
One may say this abstract view is heavy machinery for something simple; but ``we'll quite soon meet another type of algebra'' \ts{9}{1}{47:58} \ts{9}{1}{48:19}.

\begin{example}[The Lie algebra $\mathrm{End}(V)$ and the commutator as a derivation]\label{ex:tv-lie}
(After a false start with matrices: ``nonsense''.) Take $A=\mathrm{End}(V)$ for a vector space $V$ (\cref{def:te-end}); this is a vector space, not yet an algebra. Define the product
\[
  [\ ,\ ]\colon A\times A\to A,\qquad[\varphi,\psi]:=\varphi\circ\psi-\psi\circ\varphi,
\]
the \textbf{commutator}, the difference of the two compositions, quite abstractly, no matrices or matrix multiplication needed; it is again an endomorphism and it is bilinear. It is an \emph{antisymmetric} product (the bilinear map was not required to be symmetric), so $\mathrm{End}(V)$ is an algebra; one calls it a
\textbf{Lie algebra}, since the product satisfies the \textbf{Jacobi identity}
\[
  [\varphi,[\psi,\rho]]+[\psi,[\rho,\varphi]]+[\rho,[\varphi,\psi]]=0,
\]
``quickly calculated'' and following from the definition. In general any algebra whose product satisfies the Jacobi identity (and is antisymmetric) is a Lie algebra; this is one particular way to produce one \ts{9}{1}{48:19} \ts{9}{1}{49:01} \ts{9}{1}{49:42} \ts{9}{1}{50:44} \ts{9}{1}{51:25} \ts{9}{1}{52:06}.
Now the derivation: for a fixed $H\in A$ the map $D:=[H,\cdot\,]$ (a bilinear map with one slot left open is a linear map) is linear, and using the Jacobi identity one finds
\[
  D\bigl([a,b]\bigr)=[D(a),b]+[a,D(b)],
\]
which is the Leibniz rule with the product $\bullet=[\ ,\ ]$ (were the product written $\bullet$ one would see $D(a\bullet b)=D(a)\bullet b+a\bullet D(b)$). So $D=[H,\cdot\,]$ is a derivation of the Lie algebra $A$ into itself \ts{9}{1}{52:06} \ts{9}{1}{52:29} \ts{9}{1}{52:50} \ts{9}{1}{53:15} \ts{9}{1}{53:36} \ts{9}{1}{53:57}.
\end{example}
\begin{exercise}[Jacobi identity and the derivation property]\label{ex:tv-jacobi}
\emph{Not said in class.} Expand the three double commutators of \cref{ex:tv-lie} in terms of compositions, check that the twelve terms cancel in pairs, and deduce $[H,[a,b]]=[[H,a],b]+[a,[H,b]]$.
\end{exercise}
\begin{remark}[Where one has met derivations]\label{rem:tv-physics}
Where does one know this kind of derivation from? From classical mechanics: the \emph{Poisson bracket} is a derivation on the algebra of functions on phase space, the \emph{observables}; and only this is the reason why in classical mechanics one can identify the time derivative with the Poisson bracket with the Hamiltonian, since a
derivative must satisfy the product rule and so the bracket has to do so too. And from quantum mechanics: $\mathrm d/\mathrm dt$ is to be identified with $[H,\cdot\,]$ acting on observables, now self-adjoint operators. The idea of a derivation on an algebra satisfying the product rule has been met in classical mechanics, in quantum mechanics and now
in differential geometry: ``an extremely common concept, and maybe worth once and for all clarifying its algebraic status'' \ts{9}{1}{53:57} \ts{9}{1}{54:20} \ts{9}{1}{55:08} \ts{9}{1}{55:28} \ts{9}{1}{56:10} \ts{9}{1}{56:31}.
\end{remark}
\begin{coursenote}[Factors in the quantum case]
In quantum mechanics the Heisenberg equation is $\mathrm dA/\mathrm dt=\tfrac{i}{\hbar}[H,A]$ (and in classical mechanics $\mathrm df/\mathrm dt=\{f,H\}$, with the opposite order); the lecture identifies them only up to such factors and signs. Not said in class.
\end{coursenote}
The main point of the aside: the tangent vectors are derivations. Any questions so far? No \ts{9}{1}{56:31} \ts{9}{1}{56:51}. This ends the aside \ts{9}{1}{56:51}.

% ==================================================================
\section{The dimension of the tangent space}\label{sec:tv-dim}

An important theorem \ts{9}{1}{57:17}. The tangent space at another point is a vector space too, and as a vector space it has a dimension (the size of a basis, if finite).

\begin{theorem}[Dimension of the tangent space]\label{thm:tv-dim}
Let $M$ be a smooth manifold and $p\in M$. Then
\[
  \dim_{\text{vector space}}T_pM=\dim_{\text{manifold}}M.
\]
\ts{9}{1}{57:17} \ts{9}{1}{57:38} \ts{9}{1}{57:59}
\end{theorem}
These are two entirely different notions of dimension, and there is nothing obvious about the statement: the vector space dimension has technically nothing to do with the manifold dimension, which was the $\R^D$ of the charts, and the charts were not considered vector spaces, only topological spaces (\cref{ch:l06}). In principle
one could call them $\mathrm{Dim}$ and $\dim$; it needs a proof, and the strategy is to construct a vector space basis from a chart \ts{9}{1}{58:19} \ts{9}{1}{58:41} \ts{9}{1}{59:04}.

\paragraph{The chart-induced curves.} Choose a chart $(U,x)$ with $p\in U$, which one may take to be \textbf{centred at $p$}, i.e.\ $x(p)=0\in\R^D$ (without loss of generality: subtract $x(p)$ from the chart; terribly convenient). Consider $D=\dim M$ curves $\gamma_{(a)}\colon\R\to U\subseteq M$, $a=1,\dots,D$ (the index in brackets,
not to be confused with a component: it labels the curve) with the property that in the chart
\[
  x^b\bigl(\gamma_{(a)}(\lambda)\bigr)=\delta^b_a\,\lambda,\qquad b=1,\dots,D,
\]
i.e.\ in the chart the curve has all components $0$ except the $a$-th, which grows linearly with the parameter. ($\delta^b_a=1$ if $a=b$, $0$ otherwise, ``just a short notation for an array of equations''; this was asked by a student.) In $\R^D$, the Cartesian product of $D$ copies of $\R$, one has a clear notion of first, second, \dots\ component;
in the chart these curves are the coordinate axes through $0$, and $\gamma_{(a)}=x^{-1}\circ(\text{the $a$-th axis})$ are the \textbf{coordinate lines through $p$} drawn on the manifold. Another chart produces another coordinatisation. One could have drawn this picture much earlier, to get an idea of the back and forth between the charts and the manifold
\ts{9}{1}{59:27} \ts{9}{1}{1:00:09} \ts{9}{1}{1:01:11} \ts{9}{1}{1:02:35} \ts{9}{1}{1:03:38} \ts{9}{1}{1:06:26}.
Define the tangent vectors $E_a:=X_{\gamma_{(a)},p}$; all the curves pass through $p$ at $\lambda=0$ (this is why a chart centred at $p$ is used) \ts{9}{1}{1:04:40} \ts{9}{1}{1:05:02} \ts{9}{1}{1:05:22} \ts{9}{1}{1:05:43}.

Now see how the $a$-th tangent vector acts on a smooth function $f$:
\begin{align*}
E_a(f)&=(f\circ\gamma_{(a)})'(0)=\bigl((f\circ x^{-1})\circ(x\circ\gamma_{(a)})\bigr)'(0)\\
&=\partial_b\bigl(f\circ x^{-1}\bigr)\bigl(x(p)\bigr)\cdot\bigl(x^b\circ\gamma_{(a)}\bigr)'(0)=\partial_b\bigl(f\circ x^{-1}\bigr)\bigl(x(p)\bigr)\,\delta^b_a\\
&=\partial_a\bigl(f\circ x^{-1}\bigr)\bigl(x(p)\bigr),
\end{align*}
by inserting the identity $x^{-1}\circ x$, the associativity of the composition, the chain rule (the inner function is $x\circ\gamma_{(a)}\colon\R\to\R^D$, evaluated at $x(\gamma_{(a)}(0))=x(p)$) and $(x^b\circ\gamma_{(a)})'(0)=\delta^b_a$
\ts{9}{1}{1:06:47} \ts{9}{1}{1:07:28} \ts{9}{1}{1:07:49} \ts{9}{1}{1:08:30} \ts{9}{1}{1:08:51} \ts{9}{1}{1:09:33}.

\begin{definition}[New concepts: $\partial/\partial x^a$, the chart-induced tangent vectors]\label{def:tv-ddx}
For a chart $(U,x)$ with $p\in U$ and $a=1,\dots,\dim M$ the tangent vector $\bigl(\partial/\partial x^a\bigr)_p\in T_pM$ is defined by
\[
  \Bigl(\frac{\partial}{\partial x^a}\Bigr)_p(f):=\partial_a\bigl(f\circ x^{-1}\bigr)\bigl(x(p)\bigr)\qquad\text{for }f\in C^\infty(M).
\]
By the computation above $\bigl(\partial/\partial x^a\bigr)_p=X_{\gamma_{(a)},p}=E_a$ \ts{9}{1}{1:09:57} \ts{9}{1}{1:15:33} \ts{9}{1}{1:15:55}.
\end{definition}

\begin{intuition}[Two different objects: $\partial_a$ and $\partial/\partial x^a$]\label{int:tv-two-partials}
This ``at first mysterious'' remark is important and a very frequent confusion. $\partial_a$ is the $a$-th partial derivative of a function $\R^D\to\R$: it sends a smooth function $\R^D\to\R$ to a smooth function $\R^D\to\R$, ``it always did, it always will'', and it acts on functions on a Cartesian product, where it is
known which entry is the first, second, third. In physics one writes $\partial/\partial x^a$ for it. That makes no sense: it does not matter how one calls the entries when defining the function ($f(x,y,z)=x^2+y^2+z^2$ or with $b$ in place of $y$); the abstract derivative with respect to the $a$-th
entry has nothing to do with the name of the independent variable, as little as in one dimension it matters whether one writes $x$ or $\lambda$.\par
The symbol $\partial/\partial x^a$ is an \emph{entirely new symbol}, which has nothing to do with the former: it acts on functions $f$ on the manifold, $M\to\R$, for which one has no idea what the ``first entry'' would be (the manifold is a topological space); one first needs a chart map to coordinatise the manifold around $p$, obtain the chart
representative $f\circ x^{-1}$ on $\R^D$ and only then take the derivative in the $a$-th direction. So $\partial/\partial x^a$ takes a smooth function on the manifold to a real number, at the point $p$ (or the corresponding point $x(p)$); one could also produce the derivative function, but here only the value at the point is used. For the
manifold it makes sense, indeed it is necessary, that a chart has been picked. Whenever $\partial/\partial x^a$ is written in this course (and the lecturer recommends it for own calculations) the $x$ is the $x$ of the chart map: it is just a symbol, but it refers to the specific chart by which one can speak of the $a$-th component. Its value is
computed by the partial derivative of the chart representative. ``So far you never needed this; you always needed that. And here it makes sense because it refers to a chart. Very important: do not confuse the two'' \ts{9}{1}{1:09:57} \ts{9}{1}{1:10:59} \ts{9}{1}{1:12:02} \ts{9}{1}{1:13:04} \ts{9}{1}{1:14:09} \ts{9}{1}{1:15:11}.
\end{intuition}
So the tangent vectors to the coordinate-induced curves $\gamma_{(a)}$ are the chart-induced vectors $\partial/\partial x^a|_p$ \ts{9}{1}{1:15:33} \ts{9}{1}{1:15:55} \ts{9}{1}{1:16:18} \ts{9}{1}{1:16:39}. What is their meaning?

\begin{proposition}[The chart-induced vectors generate $T_pM$]\label{prop:tv-generate}
Every $X\in T_pM$ can be written as
\[
  X=X^a\,\frac{\partial}{\partial x^a}\Bigr|_p\qquad(\text{summation convention, \cref{def:te-einstein}}),
\]
with real numbers $X^a$; i.e.\ $\partial/\partial x^1|_p,\dots,\partial/\partial x^{\dim M}|_p$ is a \emph{generating system} of $T_pM$. In fact $X^a=(x^a\circ\mu)'(0)$ for any curve $\mu$ with $X=X_{\mu,p}$ \ts{9}{1}{1:16:39} \ts{9}{1}{1:16:59} \ts{9}{1}{1:17:21} \ts{9}{1}{1:17:41} \ts{9}{1}{1:23:34}.
\end{proposition}
\begin{proof}
Certainly there is a smooth curve $\mu$ through $p$ with $\mu(0)=p$ such that $X=X_{\mu,p}$: this is exactly what being an element of $T_pM$ means (there are many such curves; one is chosen) \ts{9}{1}{1:18:03} \ts{9}{1}{1:19:04} \ts{9}{1}{1:19:25} \ts{9}{1}{1:19:46}.
``How does one study a map? By letting it act on an arbitrary element of its domain'': for $f\in C^\infty(M)$,
\begin{align*}
X_{\mu,p}(f)&=(f\circ\mu)'(0)=\bigl((f\circ x^{-1})\circ(x\circ\mu)\bigr)'(0)\\
&=\partial_b\bigl(f\circ x^{-1}\bigr)\bigl(x(p)\bigr)\,(x^b\circ\mu)'(0)\\
&=(x^b\circ\mu)'(0)\,\frac{\partial}{\partial x^b}\Bigr|_p(f).
\end{align*}
The little diagram helps to see that $f\circ\mu=(f\circ x^{-1})\circ(x\circ\mu)$: restricted to the chart, $x\circ\mu$ is the representative of the curve and $f\circ x^{-1}$ that of the function, ``they come in these pairs'', and then the brackets can be dropped by associativity. The chain rule is the same calculation over and over again, and
$x(\mu(0))=x(p)$. Hence $X=X^b\,\partial/\partial x^b|_p$ with $X^b:=(x^b\circ\mu)'(0)$, a number for each $b$ \ts{9}{1}{1:19:46} \ts{9}{1}{1:20:27} \ts{9}{1}{1:21:09} \ts{9}{1}{1:21:29} \ts{9}{1}{1:22:10} \ts{9}{1}{1:22:53}.
\end{proof}
\begin{coursenote}[A slip in the board]
The board writes ``$x^b$ after $\mu$ inverse'' at one point; what is meant is $x^b\circ\mu$, not an inverse. The same remark: a few steps are written with $\gamma$ and $\mu$ interchanged. Not said in class.
\end{coursenote}
So every tangent vector is a linear combination, constructively found, of the chart-induced vectors; the numbers $X^b$ are the components with respect to this generating system (``I almost said basis, but we haven't shown that yet'') \ts{9}{1}{1:23:34} \ts{9}{1}{1:23:55}.
This is what one always did in physics: take a curve, not in the manifold abstractly but in a chart representation (the particle flies along the curve $\mu$, the first component is this function of time, the second is that function of time) and derive componentwise; the components of the tangent vector, physically the derivatives,
with respect to the chart-induced basis of the tangent space. But since one usually takes $\R^D$ as the basic space, every tangent space can be identified with $\R^D$ itself and everything is mixed up; this is the real logic behind it \ts{9}{1}{1:23:55} \ts{9}{1}{1:24:15} \ts{9}{1}{1:24:36} \ts{9}{1}{1:24:57}.
The generating system is not unique: building linear combinations of certain types gives other ones; this one is very special, induced from the coordinates, but with the advantage that derivatives of curves can be calculated in the chart picture. Later one will consider bases of the tangent space that are \emph{not} coordinate
induced, especially for vector fields on extended regions (at a point a basis can always be constructed from coordinates) \ts{9}{1}{1:25:17} \ts{9}{1}{1:25:38} \ts{9}{1}{1:25:58} \ts{9}{1}{1:26:18}.

\paragraph{Linear independence.} Suppose $\lambda^a\,\partial/\partial x^a|_p=0$ (summation over $a$; the $\lambda^a$ are real numbers, a different $\lambda$ from the curve parameter). If the vectors are linearly independent one must show that every $\lambda^a=0$. The left-hand side is an element of $T_pM$, hence an operator, so
the equation is an operator equation; it is handled by letting it act on an arbitrary element of the domain; and here one chooses very specific functions, the \emph{coordinate functions} $x^b$ \ts{9}{1}{1:26:39} \ts{9}{1}{1:27:00} \ts{9}{1}{1:27:21} \ts{9}{1}{1:27:42}.

\begin{remark}[Is it legal to act on $x^b$?]\label{rem:tv-legal}
A lawyer would say ``it depends''; the answer is yes, but not obviously. The $b$-th component $x^b\colon U\to\R$ of the chart map is continuous (the chart map is a homeomorphism onto its image, and a component of a homeomorphism to $\R^D$ is continuous; it cannot be a homeomorphism, since the dimensions differ). But $T_pM$ acts only on
\emph{smooth} functions $C^\infty(M)$, and by construction $x^b$ was only continuous on the manifold: one would not even know how to say it is differentiable. However, smoothness of a function $f$ on the manifold is \emph{decided} by looking at $f\circ x^{-1}$ in a chart: $f$ is smooth if $f\circ x^{-1}$ is. For $f=x^b$:
\[
  x^b\circ x^{-1}\colon x(U)\to\R,\qquad(a^1,\dots,a^D)\longmapsto a^b,
\]
the \emph{projection to the $b$-th component}, which is of course smooth. Since the definition of smoothness is well defined (\cref{prop:df-welldef}), the coordinate functions are smooth, and acting on them is legal; one could say that the definition of smooth function is chosen such that the chart component maps are smooth.
(This is a favourite for an oral exam: ``at first this can be very confusing''.) (A technical caveat: $x^b$ is defined on $U$ only; to act on it as an element of $C^\infty(M)$ one needs a smooth function on $M$ agreeing with it near $p$, or to let $T_pM$ act on functions defined near $p$; neither is discussed in the lecture; not said in class [?].)
\ts{9}{1}{1:27:42} \ts{9}{1}{1:28:46} \ts{9}{1}{1:29:50} \ts{9}{1}{1:30:34} \ts{9}{1}{1:31:36} \ts{9}{1}{1:32:39}.
\end{remark}
Now one can act: by the definition of the symbol, $\partial/\partial x^a|_p(x^b)=\partial_a(x^b\circ x^{-1})(x(p))$, and $x^b\circ x^{-1}$ is the projection to the $b$-th component, so the derivative in the $a$-th direction is
\[
  \frac{\partial}{\partial x^a}\Bigr|_p(x^b)=\delta^b_a.
\]
Hence $0=\lambda^a\,\partial/\partial x^a|_p(x^b)=\lambda^a\delta^b_a=\lambda^b$ for each $b$: ``you start on the assumption that this guy is zero, as an operator, so it must be zero if it acts on particular functions; we picked clever ones; we derive $\lambda^b=0$'' \ts{9}{1}{1:32:59} \ts{9}{1}{1:33:20} \ts{9}{1}{1:33:41} \ts{9}{1}{1:34:25} \ts{9}{1}{1:34:45} \ts{9}{1}{1:35:06}.
\begin{proof}[Proof of \cref{thm:tv-dim}]
$\partial/\partial x^1|_p,\dots,\partial/\partial x^{\dim M}|_p$ generate $T_pM$ (\cref{prop:tv-generate}) and are linearly independent (above); they are a basis, with $\dim M$ elements. Hence $\dim T_pM=\dim M$ \ts{9}{1}{1:35:06} \ts{9}{1}{1:35:27}.
\end{proof}
Questions? None. \ts{9}{1}{1:35:27} \ts{9}{1}{1:35:56}.
\begin{coursenote}[Students and the dots]
At \ts{9}{1}{34:44} there is a short exchange about whether something was said ``in between, like other dots''; a technical remark about the notation on the board, without content.
\end{coursenote}

% ==================================================================
\section{Components, change of chart and the myth of position vectors}\label{sec:tv-components}

\begin{definition}[New concepts: components of a tangent vector]\label{def:tv-components}
If a vector $X\in T_pM$ is written $X=X^a\,\partial/\partial x^a|_p$, the real numbers $X^1,\dots,X^{\dim M}$ are the \textbf{components} of $X$ with respect to the tangent space basis induced from the chart $(U,x)$ (where of course $p\in U$) \ts{9}{1}{1:35:56} \ts{9}{1}{1:36:17} \ts{9}{1}{1:36:39} \ts{9}{1}{1:37:03} \ts{9}{1}{1:37:23}.
\end{definition}
The problem sheet will ask for the change of components under a change of charts; a very obvious question. The key point: a chart transition map can be highly non-linear (Cartesian to polar coordinates is a highly non-linear function $\R^D\to\R^D$); but at any given point it induces a \emph{linear} transformation on the
vector space at that point, and it must, because vectors must transform linearly. If a new chart $y$ is given in terms of the old one by $y^a=y^a(x^1,\dots,x^{\dim M})$ (physicist notation), the Jacobian matrix $\partial y^a/\partial x^b$ evaluated at the point $x(p)\in\R^D$ governs the change of the basis of $T_pM$ (``kill me, I don't know whether it is the matrix $A$ or
$B$ of last time'': see \cref{prop:te-transform}); the Jacobian is a square array of functions, but it is evaluated at this one point \ts{9}{1}{1:37:43} \ts{9}{1}{1:38:04} \ts{9}{1}{1:38:47} \ts{9}{1}{1:39:07} \ts{9}{1}{1:39:50} \ts{9}{1}{1:40:10}.
\begin{exercise}[Change of components under a change of chart]\label{ex:tv-change}
Let $(U,x)$ and $(V,y)$ be charts around $p$. Show that $\dfrac{\partial}{\partial y^a}\Big|_p=\dfrac{\partial x^b}{\partial y^a}\Big|_{y(p)}\dfrac{\partial}{\partial x^b}\Big|_p$ by letting both sides act on an arbitrary $f$ (chain rule applied to $f\circ y^{-1}=(f\circ x^{-1})\circ(x\circ y^{-1})$), and find the transformation of the components $X^a$.
(Problem sheet; the long version with polar coordinates is in \cref{exam:ch09}.)
\end{exercise}

\begin{intuition}[There are no position vectors]\label{int:tv-position}
This is the full point of what is always said about there being no position vectors: the position is this. On the manifold a point is a point; what are the coordinates of Paris? There are none, unless one chooses a chart: then ``it's 7, 2'', in that chart. The chart can be changed in any diffeomorphic way ($C^\infty$ maps
$\R^D\to\R^D$), highly non-linearly. The coordinates $x^a(p)$, the components of the chart map, are certainly \emph{not} components of a vector: if they were, they would transform linearly under a change of chart, and they transform highly non-linearly. The tangent vectors to curves, on the other hand, are ``totally unimpressed'' by the non-linearity of the chart
change: at the point they only see the Jacobian matrix at that point as the linear transformation of the basis of the tangent space \ts{9}{1}{1:40:10} \ts{9}{1}{1:40:31} \ts{9}{1}{1:40:51} \ts{9}{1}{1:41:12} \ts{9}{1}{1:41:32}.
\end{intuition}
In the old books a vector is ``what transforms like this'' (with the Jacobian or its inverse, depending on the convention); but positions simply do not. Only if one restricts attention to \emph{linear} changes of coordinates does the transformation of the coordinates coincide with that of the vector components (the Jacobian of a linear map is the map itself, constant),
and then one can yield to the illusion that the coordinates constitute a vector transforming linearly. But that means one can only do Cartesian coordinates and never go to polar ones, let alone anything more interesting \ts{9}{1}{1:41:32} \ts{9}{1}{1:41:53} \ts{9}{1}{1:42:14} \ts{9}{1}{1:42:35}.
Much of special relativity is done with this confusion: the Lorentz transformations are said to act on all of position space, the $x^\mu$ ``are being transformed''; no, the $x^\mu$ can be transformed in any way one wishes, and only the derivatives, the tangent vectors at specific points, transform. If the chart transition is non-linear, the linear transformation in $T_pM$ is
not the same as that in $T_qM$ at another point $q$; it differs from point to point, from tangent space to tangent space, whereas for a linear change of chart it is everywhere the same. ``We should never abuse this by arbitrary restrictions of these functions only to save our concept of a position vector: an artificial concept, pointless, it causes confusion''
\ts{9}{1}{1:42:35} \ts{9}{1}{1:42:56} \ts{9}{1}{1:43:16} \ts{9}{1}{1:43:37} \ts{9}{1}{1:43:57}.

% ==================================================================
\section*{Key concepts}
\begin{keyconcepts}
\textbf{Retrieval questions.}
\begin{enumerate}
  \item What is the vector space $C^\infty(M)$, and which product makes it an algebra? Why is the scalar multiplication different from the product?
  \item Define the tangent vector $X_{\gamma,p}$ to a curve as an operator. Why is it the ``velocity'' of $\gamma$? What happens at double speed?
  \item Define $T_pM$. Why must one prove that the sum of two tangent vectors is again a tangent vector, and how does the curve $\sigma$ do it? Why is the chart dependence of $\sigma$ harmless?
  \item Define algebra, derivation, Lie algebra. Why is $T_pM\ni X$ a derivation, and why is $[H,\cdot\,]$ one?
  \item State and prove $\dim T_pM=\dim M$: which curves and which basis are used?
  \item What is the difference between $\partial_a$ and $\partial/\partial x^a|_p$? Why can $T_pM$ act on the coordinate functions $x^b$?
  \item How do the components of a tangent vector change under a change of chart? Why are the coordinates of a point not components of a vector?
\end{enumerate}
\textbf{Concepts.}
\begin{itemize}
  \item $C^\infty(M)$ (\cref{def:tv-cinf}); tangent vector as directional derivative operator (\cref{def:tv-vector,int:tv-married,int:tv-velocity}); the tangent space and its closure (\cref{def:tv-space,prop:tv-closure}).
  \item Algebra, derivation, Lie algebra, Poisson bracket (\cref{def:tv-algebra,def:tv-derivation,ex:tv-der-tangent,ex:tv-lie,rem:tv-physics}).
  \item $\dim T_pM=\dim M$ via the chart-induced basis (\cref{thm:tv-dim,def:tv-ddx,prop:tv-generate,rem:tv-legal}); $\partial_a$ versus $\partial/\partial x^a$ (\cref{int:tv-two-partials}).
  \item Components, change of chart, no position vectors (\cref{def:tv-components,ex:tv-change,int:tv-position}).
\end{itemize}
\end{keyconcepts}

\flushside
\end{paracol}
